\documentclass[a4paper,12pt]{article}
\usepackage[T2A]{fontenc}
\usepackage[utf8]{inputenc}
\usepackage[russian]{babel}
\usepackage{amsmath,amssymb,mathtools}
\renewcommand{\familydefault}{\sfdefault}
\usepackage{icomma}
\usepackage[a4paper,left=22mm,right=22mm,top=20mm,bottom=23mm]{geometry}
\usepackage{microtype}
\usepackage{tikz}
\usetikzlibrary{calc,arrows.meta,positioning,shapes.geometric}
\usepackage{tabularx,booktabs}
\usepackage{enumitem}
\usepackage{fancyhdr}
\usepackage{setspace}
\usepackage{xcolor}
\usepackage{needspace}
\definecolor{accent}{HTML}{246C70}
\definecolor{darkaccent}{HTML}{16464C}
\definecolor{pale}{HTML}{DAE7E6}
\definecolor{textgray}{HTML}{333A3A}
\color{textgray}
\onehalfspacing
\setlength{\parindent}{0pt}
\setlength{\parskip}{4pt}
\setlist[enumerate]{leftmargin=*,itemsep=9pt,topsep=5pt}
\setlist[itemize]{leftmargin=*,itemsep=3pt,topsep=3pt}
\pagestyle{fancy}
\fancyhf{}
\renewcommand{\headrulewidth}{0pt}
\renewcommand{\footrulewidth}{0.25pt}
\fancyfoot[L]{\scriptsize \textcolor{darkaccent}{АННА БАННОВА \; / \; 09.10.26}}
\fancyfoot[R]{\scriptsize\thepage}
\setlength{\footskip}{13mm}
\newcommand{\topic}[1]{\Needspace{5\baselineskip}\vspace{5mm}{\Large\bfseries\color{darkaccent}#1}\par\vspace{1mm}\textcolor{pale}{\rule{\linewidth}{0.7pt}}\vspace{1mm}}
\newcommand{\subtopic}[1]{\Needspace{3\baselineskip}\vspace{2mm}{\normalsize\bfseries\color{accent}#1}\par\vspace{1mm}}
\newcommand{\important}[1]{\textbf{\color{darkaccent}#1}}
\begin{document}
\thispagestyle{empty}
\begin{center}
\vspace*{37mm}
{\large\color{accent}\textbf{МАТЕМАТИКА \enspace / \enspace ЕГЭ}}\\[24mm]
{\Huge\bfseries\color{darkaccent}Чек-лист}\\[8mm]
{\LARGE\bfseries Степени и корни}\\[3mm]
{\Large Преобразование выражений\\и решение уравнений}\\[28mm]
{\normalsize\color{textgray}9 октября 2026 года}\\[8mm]
{\small Путь А \enspace\textcolor{accent}{\textbullet}\enspace 10 заданий}\\[34mm]
{\normalsize\bfseries Лёвин Артём Александрович}\\[3mm]
{\normalsize эксклюзивно для Анны Банновой}
\end{center}
\begin{tikzpicture}[remember picture,overlay]
\draw[accent!60,line width=1pt] ($(current page.south west)+(12mm,11mm)$) .. controls ($(current page.south west)+(75mm,35mm)$) and ($(current page.south east)+(-75mm,5mm)$) .. ($(current page.south east)+(-12mm,9mm)$);
\draw[pale,line width=0.45pt] ($(current page.south west)+(12mm,5mm)$) .. controls ($(current page.south west)+(83mm,25mm)$) and ($(current page.south east)+(-73mm,5mm)$) .. ($(current page.south east)+(-12mm,5mm)$);
\end{tikzpicture}
\clearpage

\topic{1. Свойства степеней}
\important{Условие для общих правил:} основание $a>0$, показатели $r,s\in\mathbb R$. Если показателем степени служит целое число, правила можно применять и к отрицательному основанию, когда выражения определены.

\renewcommand{\arraystretch}{1.25}
\begin{center}
\begin{tabularx}{\linewidth}{@{}X>{\raggedright\arraybackslash}p{68mm}@{}}
\toprule
\textbf{Действие} & \textbf{Правило} \\
\midrule
Умножение степеней с одинаковым основанием & $a^r\cdot a^s=a^{r+s}$ \\
Деление степеней с одинаковым основанием & $\displaystyle\frac{a^r}{a^s}=a^{r-s}$ \\
Возведение степени в степень & $(a^r)^s=a^{rs}$ \\
Нулевой показатель & $a^0=1$ ($a\neq0$) \\
Отрицательный показатель & $\displaystyle a^{-r}=\frac1{a^r}$ \\
Дробный показатель & $\displaystyle a^{m/n}=\sqrt[n]{a^m}$, $n\in\mathbb N$ \\
\bottomrule
\end{tabularx}
\end{center}

\subtopic{Подобные слагаемые}
Складывать коэффициенты можно при \important{полностью одинаковой буквенной части}, включая показатели степеней. Например, при $m\ne0$:
\[
\frac{7m^{30}+11m^{30}}{9m^{30}}
=\frac{18m^{30}}{9m^{30}}=2.
\]
\important{Проверка:} $7m^{30}$ и $11m^{30}$ --- подобные; $7m^{30}$ и $11m^{29}$ --- уже нет.

\subtopic{Как выбирать действие}
Сначала определите, что мешает вычислению: скобки, разные показатели или корни. Затем приведите запись к степеням с одинаковым основанием, выполните преобразования и только после этого подставляйте числовые значения.

\topic{2. Дробные показатели и корни}
\subtopic{Дробь в показателе: внимательно к скобкам}
При $n>0$:
\[
\frac{6n^{1/3}}{n^{1/12}\cdot n^{1/4}}
=6n^{\frac13-\left(\frac1{12}+\frac14\right)}
=6n^{\frac4{12}-\frac1{12}-\frac3{12}}=6.
\]
\important{Важно:} знаменатель дроби даёт \emph{вычитание всей суммы} показателей. После раскрытия скобок оба знака меняются.

\clearpage
\topic{3. Переход от корней к степеням}
При $a>0$ верно $\sqrt[k]{a}=a^{1/k}$. Например:
\[
\frac{\sqrt[9]{a}\cdot\sqrt[18]{a}}{a\sqrt[6]{a}}
=a^{\frac19+\frac1{18}-\left(1+\frac16\right)}
=a^{-1}=\frac1a.
\]
Если $a=\frac12$, то значение выражения равно $\left(\frac12\right)^{-1}=2$.

\subtopic{Три обязательные проверки записи}
\begin{enumerate}[label=\arabic*)]
\item \important{Знаки:} при вычитании суммы $u-(v+w)$ получаем $u-v-w$.
\item \important{Вычисления:} при сложении показателей находите общий знаменатель.
\item \important{Скобки:} при подстановке $a=\frac12$ пишите $\left(\frac12\right)^{-1}$, чтобы степень относилась ко всему числу.
\end{enumerate}

\topic{4. Аргумент функции и подстановка}
Обозначение $g(t)$ означает значение функции \emph{в точке} $t$. В частности, если $g(t)=8^t$, то $g(x-9)=8^{x-9}$, а $g(x-11)=8^{x-11}$.
\[
\frac{g(x-9)}{g(x-11)}
=\frac{8^{x-9}}{8^{x-11}}
=8^{\,(x-9)-(x-11)}=8^2=64.
\]
\important{Типичная ошибка:} терять скобки при вычитании показателя знаменателя. Запись $x-9-x-11$ неверна.

\topic{5. Показательные уравнения}
При $a>0$ и $a\ne1$:
\[
 a^{F(x)}=a^{G(x)}\quad\Longleftrightarrow\quad F(x)=G(x).
\]
Нужно представить обе части как \important{одну степень с одним и тем же основанием}.
\begin{align*}
2^x=\frac12=2^{-1}&\quad\Rightarrow\quad x=-1,\\
2^x=\sqrt[3]{2}=2^{1/3}&\quad\Rightarrow\quad x=\frac13,\\
4\cdot2^x=8\;\Longleftrightarrow\;2^{2+x}=2^3
&\quad\Rightarrow\quad x=1.
\end{align*}
\important{Ограничение:} если основания нельзя привести к одному числу, потребуются другие методы.

\clearpage
\thispagestyle{plain}
\begin{center}
{\Large\bfseries\color{darkaccent}Алгоритм: показательное уравнение}\\[3mm]
\textcolor{accent}{От записи степеней к равенству показателей}
\end{center}
\vspace{8mm}
\begin{center}
\begin{tikzpicture}[
  >=Stealth,
  node distance=11mm,
  term/.style={ellipse,draw=accent,line width=0.85pt,fill=pale!23,align=center,text width=74mm,minimum height=14mm,inner sep=4mm,font=\small},
  action/.style={rectangle,rounded corners=2pt,draw=accent!75,line width=0.8pt,align=center,text width=78mm,minimum height=16mm,inner sep=4mm,font=\small},
  decision/.style={diamond,aspect=2.15,draw=darkaccent,line width=0.8pt,align=center,text width=40mm,inner sep=2mm,font=\small},
  side/.style={rectangle,rounded corners=2pt,draw=accent!75,align=center,text width=41mm,inner sep=3mm,font=\small},
  arrow/.style={->,thick,draw=accent!95}
]
\node (start) [term] {$a^{F(x)}=B$\\$a>0$, $a\neq1$};
\node (conv) [action,below=of start] {Представьте $B$ как степень основания $a$:\\число, дробь, корень, произведение};
\node (test) [decision,below=17mm of conv] {Получено\\$B=a^{G(x)}$?};
\node (yes) [action,below=19mm of test] {Приравняйте показатели:\\$F(x)=G(x)$};
\node (last) [term,below=of yes] {Решите уравнение и проверьте ответ};
\node (no) [side,right=16mm of test] {Этот метод пока не завершает решение. Потребуются другие инструменты.};
\draw[arrow] (start)--(conv);
\draw[arrow] (conv)--(test);
\draw[arrow] (test)--node[right,font=\small]{да}(yes);
\draw[arrow] (yes)--(last);
\draw[arrow] (test.east)--node[above,font=\small]{нет}(no.west);
\end{tikzpicture}
\end{center}
\clearpage

\topic{6. Когда неизвестное стоит в основании}
В уравнении $x^{4/3}=8$ нужно учитывать, что \important{основание $x$ может быть отрицательным}. Поскольку знаменатель показателя равен $3$, кубический корень определён для всех действительных $x$.

Положим $t=\sqrt[3]{x}$. Тогда:
\[
 t^4=8\quad\Longrightarrow\quad
 t=\pm\sqrt[4]{8},\qquad
 x=t^3=\boxed{\pm8^{3/4}}.
\]
Получены \important{два действительных решения}. Прямое возведение обеих частей в степень $\frac34$ без учёта знака могло бы потерять отрицательный корень: $\bigl(x^{4/3}\bigr)^{3/4}=|x|$.

\subtopic{Отрицательный целый показатель}
При $x\ne0$ уравнение $x^{-2}=4$ равносильно:
\[
\frac1{x^2}=4\quad\Longleftrightarrow\quad
x^2=\frac14\quad\Longleftrightarrow\quad
\boxed{x=\pm\frac12}.
\]
Знак «минус» в показателе означает \important{обратную величину}; он не означает, что основание обязано быть отрицательным.

\subtopic{Мини-памятка}
\begin{itemize}
\item Для $a^x=a^y$ можно приравнивать показатели лишь при $a>0$, $a\ne1$.
\item При сокращении степеней с буквами проверьте, что исходный знаменатель не равен нулю.
\item В уравнениях $x^{2k}=c$ при $c>0$ обычно два корня: положительный и отрицательный.
\item После подстановки найденных корней проверьте исходное уравнение.
\end{itemize}

\clearpage
\begin{center}
{\LARGE\bfseries\color{darkaccent}Тренировочный блок}\\[3mm]
{\normalsize Путь А \quad\textcolor{accent}{\textbullet}\quad 10 заданий}
\end{center}
\vspace{4mm}
\textbf{Порядок работы.} Решайте по порядку. Записывайте преобразования и итоговый ответ. В заданиях 9--10 укажите \emph{все действительные корни}.
\begin{enumerate}[label=\textbf{\arabic*.}]
\item Упростите выражение ($t\neq0$):
\[ \frac{5t^{12}+7t^{12}}{3t^{12}}. \]
\item Упростите выражение ($p>0$):
\[ \frac{12p^{5/6}}{p^{1/3}\cdot p^{1/2}}. \]
\item Найдите значение выражения при $a=9$:
\[ \frac{\sqrt[4]{a^3}\cdot\sqrt{a}}{\sqrt[4]{a}}. \]
\item Найдите значение выражения при $a=\frac14$:
\[ \frac{\sqrt[3]{a}\cdot\sqrt[6]{a}}{a^{3/2}}. \]
\item Если $f(t)=3^t$, найдите значение
\[ \frac{f(x+5)}{f(x+2)}. \]
\item Решите уравнение:
\[ 2^x=\frac{1}{\sqrt[3]{4}}. \]
\item Решите уравнение:
\[ 9\cdot3^{x-1}=27. \]
\item Решите уравнение:
\[ 5^{2x-1}=\frac{125}{\sqrt5}. \]
\item Решите уравнение в действительных числах:
\[ x^{4/3}=16. \]
\item Решите уравнение в действительных числах:
\[ x^{-4}=\frac1{16}. \]
\end{enumerate}
\clearpage
\begin{center}
{\LARGE\bfseries\color{darkaccent}Ответы для самопроверки}\\[4mm]
\end{center}
\renewcommand{\arraystretch}{1.6}
\begin{center}
\begin{tabularx}{\linewidth}{@{}>{\bfseries}p{15mm}X@{}}
\toprule
№ & Ответ \\
\midrule
1 & $4$ \\
2 & $12$ \\
3 & $9$ \\
4 & $4$ \\
5 & $27$ \\
6 & $x=-\dfrac23$ \\
7 & $x=2$ \\
8 & $x=\dfrac74$ \\
9 & $x=-8$; $x=8$ \\
10 & $x=-2$; $x=2$ \\
\bottomrule
\end{tabularx}
\end{center}
\vspace{9mm}
\important{Самопроверка:} для каждого задания оцените не только ответ, но и сохранение области определения, правильность скобок и знаков. В заданиях 9 и 10 подстановка обоих значений в исходное уравнение обязательна.
\vfill
\begin{center}
\small\color{darkaccent}Лёвин Артём Александрович эксклюзивно для Анны Банновой
\end{center}
\end{document}
